% ===============================
% Cellbox: SIGCOMM / NSDI LaTeX Draft
% ===============================

\documentclass[10pt,conference]{IEEEtran}

\usepackage{times}
\usepackage{amsmath,amssymb}
\usepackage{graphicx}
\usepackage{booktabs}
\usepackage{multirow}
\usepackage{url}
\usepackage{xcolor}

\title{Cellbox: A Cell-Based Hierarchical Scheduling and Flow-Control Fabric for Lossless Datacenter Networks}

\author{\IEEEauthorblockN{Parham Soltani}
\IEEEauthorblockA{Department of Electrical Engineering\\
Iran University of Science and Technology\\
Email: \texttt{parham.soltani@...}}}

\begin{document}
\maketitle

% ===============================
\begin{abstract}
Modern datacenter networks must simultaneously provide ultra-low tail latency, high throughput, strict fairness, and lossless operation for heterogeneous workloads such as RDMA, AI training, and RPC traffic. Existing solutions address these goals in isolation: hierarchical schedulers approximate fairness but ignore flow control, while lossless fabrics regulate admission without enforcing precise service ordering. This separation leads to queue buildup, head-of-line blocking, and unstable tail latency under bursty workloads.

We propose \emph{Cellbox}, a hardware-native, lossless datacenter switch fabric that unifies scheduling and flow control at cell granularity. Cellbox segments packets into fixed-size cells and schedules them through a hierarchical calendar-queue structure that is explicitly credit-aware. Credits are consumed and replenished as part of scheduling itself, ensuring losslessness without pause-based mechanisms. Analytical modeling shows bounded delay and queue stability, while simulations demonstrate up to 40\% lower tail latency and 2$\times$ improved fairness compared to Gearbox. An FPGA prototype demonstrates feasibility at 400~Gbps per port.
\end{abstract}

% ===============================
\section{Introduction}

Datacenter networks increasingly support workloads with stringent performance requirements, including distributed AI training, RDMA-based storage, and latency-sensitive microservices. These workloads demand lossless transport, predictable bandwidth allocation, and tight tail-latency bounds—even under extreme burstiness and incast.

Existing mechanisms decompose the problem: schedulers determine service order without downstream awareness, while lossless transport mechanisms regulate admission without enforcing fairness. At 400--800~Gbps link rates, packet-level scheduling becomes too coarse to satisfy both goals simultaneously.

This paper presents \textbf{Cellbox}, a lossless datacenter switch fabric that co-designs scheduling and flow control at cell granularity. Cellbox is a fixed-function hardware architecture; programmability is limited to slow-timescale configuration and is not required for correctness or performance.

% ===============================
\section{Background and Motivation}

\subsection{Hierarchical Scheduling}
Gearbox introduced scalable hierarchical scheduling with bounded departure-time discrepancy. However, it schedules packets without knowledge of downstream availability and relies on external mechanisms such as PFC for losslessness.

\subsection{Lossless Transport}
PFC and credit-based schemes prevent drops but operate independently of scheduling, leading to head-of-line blocking and congestion spreading.

\subsection{Opportunity}
Losslessness is fundamentally a scheduling problem. If admission and service are jointly controlled, packet drops can be eliminated without pauses.

% ===============================
\section{Cellbox Architecture}

\subsection{Design Overview}
Cellbox segments packets into fixed-size cells and schedules them through a hierarchical, credit-aware calendar queue. Cells are admitted only when downstream credits are available.

\subsection{Hierarchical Scheduling}
Scheduling is organized into per-flow, per-class, and per-port levels. Each level uses a timing wheel indexed by virtual departure time, enabling O(1) enqueue and dequeue.

\subsection{Credit-Aware Admission}
For cell $k$ of flow $i$:
\begin{equation}
T_i^k = \max(V(t), T_i^{k-1}) + \frac{C}{w_i}
\end{equation}
Cells are scheduled only if sufficient credits exist, guaranteeing losslessness.

% ===============================
\section{Analysis}

We model queue evolution as:
\begin{equation}
\frac{dQ_i}{dt} = \lambda_i - \mu_i(t)
\end{equation}
where service rate is bounded by both scheduling weight and available credits. Using Lyapunov arguments, queues are stable under admissible traffic.

% ===============================
\section{Evaluation}

We evaluate Cellbox using ns-3 simulations and an FPGA prototype.

\subsection{Simulation}
We use Clos topologies with RDMA and RPC workloads. Baselines include Gearbox, PFC+DCTCP, ExpressPass, and BFC.

\subsection{Prototype}
An FPGA-based switch fabric demonstrates feasibility at 400~Gbps per port.

% ===============================
\section{Related Work}

\subsection{Hierarchical Scheduling}
Gearbox provides scalable fairness but assumes infinite buffering. Cellbox extends hierarchical scheduling with cell granularity and integrated flow control.

\subsection{Lossless Fabrics}
PFC and ExpressPass ensure zero loss but do not enforce service ordering or fairness. Cellbox unifies these concerns in hardware.

% ===============================
\section{Why PFC Is Fundamentally Inferior}

PFC is reactive, coarse-grained, and decoupled from scheduling. It causes head-of-line blocking and congestion spreading, and provides no service guarantees. Cellbox enforces losslessness proactively at scheduling time.

% ===============================
\section{Conclusion}

Cellbox demonstrates that fairness, bounded delay, and losslessness can be achieved through a unified, cell-based scheduling fabric. By embedding credit awareness directly into hierarchical scheduling, Cellbox eliminates pause-based losslessness while remaining hardware-efficient.

\bibliographystyle{IEEEtran}
\bibliography{cellbox}

\end{document}
